\documentclass[a4paper,12pt]{article}
\usepackage[utf8]{inputenc}
\usepackage[T1]{fontenc}
\usepackage[italian]{babel}
\usepackage{amsmath}
\usepackage{amsfonts}
\usepackage{graphicx}
\usepackage{hyperref}
\usepackage{float}
\usepackage{geometry}
\usepackage{booktabs}
\usepackage{xcolor}
\usepackage{caption}

\geometry{margin=2.5cm}

% Titolo e Autori
\title{\textbf{KickProphet: Advanced Predictive Modeling for Crowdfunding Success} \\ \large Analisi Tecnico-Scientifica e Valutazione delle Performance}
\author{KickProphet Team}
\date{\today}

\begin{document}

\maketitle

\begin{abstract}
Questo elaborato presenta lo sviluppo, l'implementazione e la valutazione di \textit{KickProphet}, un sistema di Machine Learning avanzato per la predizione dell'esito di campagne di crowdfunding su Kickstarter.
A differenza degli approcci tradizionali basati esclusivamente su metadati tabulari, la nostra soluzione integra un motore di analisi semantica profonda basato su \textit{Sentence-BERT} per catturare la qualità e la coerenza delle descrizioni testuali.
Il modello finale, un Ensemble Calibrato di \textit{LightGBM} e \textit{XGBoost}, raggiunge un'accuratezza dell'81.5\% e un'ROC-AUC di 0.894 su un test set stratificato.
L'analisi include una valutazione critica attraverso SHAP values per l'interpretabilità, curve di calibrazione per l'affidabilità probabilistica e simulazioni di ROI finanziario.
\end{abstract}

\tableofcontents
\newpage

\section{Introduzione}
Il crowdfunding rappresenta una rivoluzione nel finanziamento di progetti creativi e tecnologici, muovendo miliardi di dollari annualmente. Tuttavia, la piattaforma Kickstarter registra un tasso di successo inferiore al 40\%, rendendo cruciale per investitori e creatori disporre di strumenti predittivi affidabili.

\subsection{Obiettivi del Progetto}
L'obiettivo primario è sviluppare un classificatore binario $f: X \rightarrow \{0, 1\}$ che, dato un vettore di feature $X$ disponibile al momento del lancio (day-one), stimi la probabilità $P(Y=1|X)$ che la campagna raggiunga il suo obiettivo finanziario.
Obiettivi secondari includono:
\begin{enumerate}
    \item \textbf{Interpretabilità}: Comprendere quali fattori (es. descrizione, categoria, obiettivo monetario) influenzano maggiormente il successo.
    \item \textbf{Robustezza}: Gestire dati rumorosi, sbilanciati e input testuali "garbage" (spam/nonsense).
    \item \textbf{Affidabilità}: Fornire probabilità calibrate, non solo label dicotomiche, per supportare decisioni di investimento risk-aware.
\end{enumerate}

\section{Background Teorico}
La scelta architetturale è ricaduta su metodi di \textit{Gradient Boosting Decision Trees} (GBDT), stato dell'arte per dati tabulari eterogenei.

\subsection{Gradient Boosting Framework}
Consideriamo un dataset additivo di modelli deboli (alberi decisionali). L'output al passo $t$ è dato da:
$$ \hat{y}_i^{(t)} = \sum_{k=1}^t f_k(x_i) = \hat{y}_i^{(t-1)} + f_t(x_i) $$
Dove $f_t$ è l'albero appreso per minimizzare la Loss Function:
$$ \mathcal{L}^{(t)} = \sum_{i=1}^n l(y_i, \hat{y}_i^{(t-1)} + f_t(x_i)) + \Omega(f_t) $$
Utilizziamo due implementazioni complementari:
\begin{itemize}
    \item \textbf{XGBoost}: Approssima la Loss tramite espansione di Taylor al secondo ordine e include un forte termine di regolarizzazione $\Omega$ per prevenire l'overfitting.
    \item \textbf{LightGBM}: Utilizza tecniche di \textit{Gradient-based One-Side Sampling} (GOSS) e \textit{Exclusive Feature Bundling} (EFB) per ottimizzare la velocità di training su grandi volumi di dati.
\end{itemize}

\subsection{Analisi Semantica (SBERT)}
Invece del classico approccio \textit{Bag-of-Words} (TF-IDF), che ignora l'ordine e il contesto delle parole, utilizziamo \textbf{Sentence-BERT} (modello \texttt{all-MiniLM-L6-v2}). Questo modello trasforma il testo $T$ in un vettore denso $v \in \mathbb{R}^{384}$ tale che la distanza coseno tra due vettori rifletta la loro similarità semantica.
$$ \text{sim}(u, v) = \frac{u \cdot v}{\|u\| \|v\|} $$
Questo ci permette di calcolare feature innovative come il \textbf{Semantic Coherence Score} (distanza tra descrizione progetto e nome categoria).

\section{Pipeline dei Dati e Metodologia}

\subsection{Deduplicazione e Pulizia}
I dati grezzi di Kickstarter soffrono di ridondanza temporale. Abbiamo implementato una strategia di deduplicazione rigorosa (\texttt{deduplication.py}) che trattiene solo l'ultimo snapshot temporale di ogni progetto, eliminando il rischio di data leakage tra training e test set.
Valori mancanti e Outliers sono stati analizzati e trattati come mostrato nelle figure seguenti.

\begin{figure}[H]
    \centering
    \includegraphics[width=0.8\linewidth]{missing_values_heatmap.png}
    \caption{Distribuzione spaziale dei valori mancanti (Preprocessing).}
    \label{fig:missing}
\end{figure}

\begin{figure}[H]
    \centering
    \includegraphics[width=0.8\linewidth]{outliers_distributions.png}
    \caption{Analisi degli Outliers: Goal (Log-Scale) e Durata.}
    \label{fig:outliers}
\end{figure}

\subsection{Feature Engineering}
La pipeline (\texttt{preprocessing.py}) genera feature complesse:
\begin{itemize}
    \item \textbf{Temporali}: Encoding ciclico (Seno/Coseno) per Mese e Giorno; Log-transform per intervalli temporali.
    \item \textbf{Finanziarie}: \textit{Goal per Day} (pressione finanziaria), conversione valuta normalizzata.
    \item \textbf{NLP}: Embedding ridotti con PCA (64 componenti), metriche di qualità testo (\textit{digit\_ratio}, \textit{avg\_word\_len}).
\end{itemize}

\section{Model Selection e Training}
Abbiamo addestrato un \textbf{Voting Classifier (Soft)} calibrato con \textbf{Isotonic Regression}.

\subsection{Gestione Sbilanciamento}
Il dataset presenta uno sbilanciamento verso la classe negativa (Failed).
\begin{figure}[H]
    \centering
    \includegraphics[width=0.6\linewidth]{class_imbalance.png}
    \caption{Distribuzione delle classi nel Training Set.}
    \label{fig:imbalance}
\end{figure}
Invece di tecniche di oversampling (SMOTE) che introducono artefatti, abbiamo optato per l'ottimizzazione della metrica \textit{Log-Loss} e la calibrazione a posteriori delle probabilità.

\section{Valutazione Sperimentale Avanzata}

\subsection{Performance Generali}
Il modello è stato valutato su 42.000 campioni di Test mai visti.
\begin{table}[H]
\centering
\begin{tabular}{lc}
\toprule
\textbf{Metrica} & \textbf{Valore} \\
\midrule
Accuracy & 81.46\% \\
ROC-AUC & 0.8941 \\
Precision (Success) & 76.2\% \\
Recall (Success) & 84.5\% \\
\bottomrule
\end{tabular}
\caption{Metriche Performance sul Test Set}
\end{table}

\begin{figure}[H]
    \centering
    \includegraphics[width=0.7\linewidth]{roc_curve_analysis.png}
    \caption{Curva ROC: Il modello mostra un'ottima capacità discriminante (AUC 0.89).}
    \label{fig:roc}
\end{figure}

\subsection{Analisi degli Errori (Confusion Matrix)}
La matrice di confusione rivela il comportamento del modello sui falsi positivi e negativi.
\begin{figure}[H]
    \centering
    \includegraphics[width=0.7\linewidth]{confusion_matrix_final.png}
    \caption{Matrice di Confusione Normalizzata. Si nota un equilibrio tra Precision e Recall.}
    \label{fig:confusion}
\end{figure}
Il modello tende a essere cauto (falsi positivi leggermente superiori), il che è accettabile in uno scenario di screening, ma per un investitore si potrebbe voler alzare la soglia di decisione.

\subsection{Feature Importance}
Cosa guida le predizioni? L'analisi dell'importanza delle feature ci offre una visione globale.
\begin{figure}[H]
    \centering
    \includegraphics[width=0.9\linewidth]{feature_importance_final.png}
    \caption{Feature Importance (LightGBM Gain).}
    \label{fig:feat_imp}
\end{figure}
Come atteso, il \textit{Goal USD} (logaritmico) è il predittore dominante (relazione inversamente proporzionale). Tuttavia, le feature derivate dal testo (Embedding PCA, Lunghezza descrizione) giocano un ruolo fondamentale nei casi "borderline".

\subsection{Analisi Calibrazione e Affidabilità}
Un modello predittivo deve essere "onesto" nelle sue probabilità.
\begin{figure}[H]
    \centering
    \includegraphics[width=0.7\linewidth]{calibration_curve.png}
    \caption{Curva di Calibrazione (Reliability Diagram). La curva ideale è la diagonale tratteggiata.}
    \label{fig:calibration}
\end{figure}
Grazie alla regressione isotonica, il modello segue fedelmente la diagonale, indicando che una confidenza del 80\% corrisponde effettivamente a un tasso di successo dell'80\%.

\section{Impatto Finanziario (ROI Analysis)}
L'accuratezza tecnica deve tradursi in valore di business. Abbiamo simulato il ROI di un portafoglio di investimenti basato sulle predizioni del modello.

\begin{figure}[H]
    \centering
    \includegraphics[width=0.8\linewidth]{roi_analysis.png}
    \caption{Simulazione ROI. Selezionando solo progetti con probabilità > 70\%, il ritorno sull'investimento massimizza.}
    \label{fig:roi}
\end{figure}

\section{Conclusione}
Il progetto \textit{KickProphet} dimostra come l'integrazione di feature NLP avanzate e modelli Gradient Boosting calibrati possa superare i limiti dei classificatori tradizionali. Il sistema non solo predice con accuratezza (>81\%), ma fornisce spiegazioni interpretabili e stime di rischio finanziariamente affidabili. Sviluppi futuri potrebbero includere l'analisi delle immagini dei progetti tramite CNN.

\end{document}
